\scanpage{9}
rer que pour $d$ assez grand, ce nombre est non nul (ce qui impliquera \deleted{\unsure{al}}
alors\deleted{,} -- et est sans doute équivalent -- que le motif $E(Y)$ pour tous
ces $Y$ est de \uline{niveau} maximum $n-1$, ou encore que le motif $E(Y)$ pour $Y$
générique est de niveau $n-1$). Sous la forme de (8.5.6)\deleted{\unsure{*}}, on voit que
ceci est assuré dès que l'on sait que
\[
\dim E(Y,\underline O_Y)_{\mathrm{SS}}\deleted{\unsure{\longrightarrow+\infty}}\inserted{>0}\qquad\text{pour}\quad d\longrightarrow+\infty,
\]
où le dim est pris bien entendu avec sa valeur générique.
\deleted{\unsure{On sait en effet} (cf n° \unsure{7}) que $\dim\unsure{E(Y,\underline O_Y)_{\mathrm{ss}}}>\unsure{\rho|0}$ pour $d$ grand, et que $\varphi\ne0$ \unsure{pour $d$ grand} (avec \unsure{$E(Y)$} \unsure{de niveau} $\unsure{n}-1$) dès que l'on sait que $\dim(\operatorname{Tors}^{\unsure{n+1|n-1}}(X)\unsure{\otimes}k)_{\mathrm{ss}}<\unsure{\rho|p}-1$.}
\margin{Or cela est établi au n° 7}

8.5.9. Le corps de base $k$ étant à nouveau quelconque (même pas nécessai-
rement parfait), désignons, pour un degré $d$ donné, par $U_d$ l'ouvert de
l'espace projectif des hypersurfaces de degré $d$ dans $P^r_k$, formé des
hypersurfaces dont l'intersection avec $X$ est \deleted{non singulière}\inserted{lisse} de \unsure{dimension}
$n-1$, et telles que $E(Y,\underline O_Y)_{\mathrm{SS}}$ soit de rang maximum. Alors les suites
exactes (8.5.1), pour les $Y$ provenant de points de $U_d$, proviennent d'u\unsure{ne}
suite exacte de Modules localement libres \inserted{sur $U_d$} à opération $p$-linéaire semi-
simple, ou de façon équivalente (cf 8.4.2) à une suite exacte de systè-
mes locaux \inserted{de vectoriels sur $\underline F_p$} sur $U_d$ pour la topologie étale, \deleted{formés} ou enfin, une fois
choisi un point géométrique de $U_d$ (par exemple via une clôture \unsure{séparable}
du point générique) une suite exacte de représentations \inserted{de $G=\pi_1(U_d)$} dans des vectori-
els de dimension finie sur $\underline F_p$. \deleted{On} On trouve un homomorphisme \inserted{$\varphi$} du
système local défini par les $E(Y)\otimes_W k'$ \deleted{dans}\inserted{\uline{sur}} celui défini par les
$E(Y,\underline O_Y)$, dont les noyau \deleted{et le conoyau sont des}\inserted{est un} système\deleted{s} loca\deleted{ux}\inserted{l} \uline{\unsure{constant}}
(i.e. provenant du corps de base $k$). Lorsque $k$ est fini, on peut suppo-
ser, quitte à passer à une extension finie de $k$, que les valeurs propres
du frobenius relatif à $k$ opérant sur ce\deleted{s} dernier\deleted{s} \deleted{\unsure{deux}} sont toutes
égales à $1$ (il suffit que les valeurs propres de frobenius relatif à $k$
opérant sur \deleted{les} $H^{n\deleted{\unsure{-1}}}(X,\underline O_X)_{\mathrm{SS}}$ \deleted{et sur les deux membres de (8.5.3)} soient
\mspage{8}
